% !TeX root = ../papers/nc-p3.tex
\section{頂点上の導手と母線上での核座標の復元}
\label{sec:cofactor-eighteen-cone-pencil}
\begin{lemma}\label{lem:cofactor-cone-pencil-collision}
二次錐の母線を $t\in\mathbb P^1$ で表す。一般核点の二重被覆を局所的に
$w^2=P(t,z)$、その核点を母線上のアフィン座標 $z$ で表し、入力四座標が
\[
                A(t)+zB(t)+w C(t)
\]
であるとする。$P$ は $z$ について高々二次、または $z$ に依存しない。
同次化した各母線上で入力は有限、総入力次数二とする。
$B$ が固定入力直線 $\Lambda$ への次数一のパラメータ表示であり、
非零の $C$ もその同じ二次元ベクトル空間に値を持つなら矛盾する。
\end{lemma}
\begin{proof}
共通零を除いた $[C]$ は底から $\Lambda$ への射である。
$[B]$ と一致する点 $t_0$ がある（対角因子を引き戻すと次数 $1+\deg[C]$。
恒等的に一致する場合は一般点を取る）。局所パラメータ $s=t-t_0$ として
$C=s^m(C_0+O(s))$, $C_0=\lambda B(t_0)$, $\lambda\ne0$。
この母線の入力像は直線でなければならない。一点なら有限性に反する。
その一般点の入力階数は三である（階数二曲線は整で次数十八）。
従って近傍で核点は入力だけの正則関数であり、錐の頂点を避ければ $z$ も正則に復元される。

$m=0$ で $P(t_0,z)$ が恒等零でなければ、直線の座標は $z+\lambda w$ である。
この直線から核母線への写像は次数一になる。実際被覆母線から両方の直線への総次数は二であり、
可約なら各成分の次数を足しても同じである。従って一般の同じ $z$ を持つ二点が
異なる $z+\lambda w$ を持つことと両立しない。

$m\ge1$ で同じく一般 $z$ で非分岐なら、$s$ 進の二枝を取る（剰余体を拡大してよい）。
第一枝の値を $w$、他方の枝では $z'=z+2\lambda s^m w$ としその平方根を
$w'=-w+O(s^m)$ に取る。両入力は $s^{m+1}$ を法として等しいが $z,z'$ は異なる。
垂直分岐の場合は分岐式の被約性から一般 $z$ で $s$ の位数が一である。
$s=v^2$, $w=v W(v,z)$ として $W$ は単元の平方根を選ぶ。
同じ $z'=z+2\lambda s^m w$ と反対の平方根を使えば、
入力は $v^{2m+2}$ を法として等しく $z,z'$ は異なる。これは $m=0$ にも適用できる。
いずれも階数三での正則復元に反する。
\end{proof}

\begin{theorem}\label{thm:cofactor-eighteen-no-quadric}
$d=18$ では二次錐も核像にならない。従って十八次の曲面核像に残るのは
非正規三次または四次で、入力逆写像がその特異曲線の一般点の少なくとも一つで降りない場合だけである。
三次の場合は核切断の算術種数が三である。
\end{theorem}
\begin{proof}[一般母線の連結二重被覆]
定理\ref{thm:cofactor-eighteen-cone-ruling-degree}の $T\to\mathbb F_2$ を使う。
$L=g^*(E+2F)$、$U=g^*F$、$H\equiv L+U$、入力 $f$ は有限である。
$a=1$ なら $T$ は有理で有理二重点のみ（先の二重被覆列による $\chi=1$ の計算）。
従って数値同値は解消のPicard群で線形同値になり、$H=L+U$ が $T$ 上でも成立する。
$\mathbb F_2$ の母線座標を $[r:z]$ とし $r=0$ を負切断とする。
核点は $[r u^2:r uv:r v^2:z]$、$t=[u:v]$ である。跡分解により
$r=1$ の局所表示で入力は
\[
        A_3(t)+zB_1(t)+w C_{3-b}(t),\qquad w^2=P(t,z),
\]
添字は二項式の次数で、$w$ は母線に沿って一次の重み、$P$ は相対二次である。
$C$ が零なら入力は二重被覆を経由するので不可能である。

対称余核の跡の有効因子は $g^*(E+(2-b)F)$ の切断である。
負切断が分岐に入る場合、その上の被約切断 $s$ でこの因子は消える。
実際跡分解で反不変な切断は $\mathcal O((2-2b)F)$ の切断に $w$ を掛けたもの、
不変切断は $|E+(2-b)F|$ の切断であり、いずれも $s$ 上で零である。
$s$ の入力次数一なので $B_1(t)$ は固定直線 $\Lambda$ の次数一の表示である。
一般点は階数三で、核頂点を持つ他の正規化点は同じ入力の上にない。
$C(t)$ が $\Lambda$ のベクトル空間の外に出るなら、$s$ からの法方向には $w$ が局所パラメータで、
入力の微分は $C(t)$ 方向に出る（$r$ は二次以上、ここでは $z=1$ のチャートを使う）。
よって有限入力射はその一般点で一対一かつ非分岐になり同型である。
そこで跡が消えることに反する。従って $C(t)$ も $\Lambda$ 内で、補題が使える。

負切断が分岐に入らない場合は $b=2$ で、分岐と $E$ は互いに素である。
$g^{-1}E$ は二本の互いに素な切断であり、跡因子は正確にその和である
（$g^*E$ の切断は不変なただ一つの切断だけ）。
入力は二切断を各々直線へ同型に写す。その二直線が違えば一般の各像で階数三、
ただ一つの逆像を持ち、核座標は負切断の法方向に微分非零である。
従って入力接方向と核の法方向で正規化の微分を分離し、階数三ゆえ入力射自体がそこで同型となり矛盾。
両直線は同じ $\Lambda$ である。負切断上では入力は $B_1(t)\pm c C_1(t)$
（$c\ne0$ は定数）なので $B,C$ とも $\Lambda$ に入る。
$B$ が零を持つ場合もここでは除ける。実際その $t_0$ で $B\pm cC$ は同じ射影点を与える。
より直接、$B$ が退化した一次列である場合を次段落でまとめて扱う。

$B,C$ を同じ二次元空間の一次列とすると二次式 $\det(B,C)$ に零点がある。
$B\ne0$ の点で、共通因子を除いた $[C]$ と $[B]$ が一致すれば補題の局所証明が使える。
$B$ が基点を持つ場合はそこに直接矛盾がある。その母線に垂直分岐はない
（垂直分岐なら $E$ と交わる）。$B(t_0)=0$ なら入力直線座標は $w$ の一次式、
核座標は $z$ である。有限性からいずれも総次数二なので一般入力点での復元は両者を
一次分数変換で結ぶ。しかし一般の同じ $z$ の二平方根は異なる入力点を与える。
$A(t_0),C(t_0)$ が比例なら有限性に反する。従って $B$ には基点がなく補題を適用できる。
\end{proof}

\begin{proof}[一般母線が二本に分かれる被覆]
$a=0$ の場合、$T$ は底の二重被覆 $B_0\to\mathbb P^1$ を用いた
$\mathbb F_2$ の底変換である。$B_0$ は滑らかで種数 $g_0=b-1$。
負切断の引戻し $s$ は $B_0$ 自身で $s^2=-4$、$T$ 上の有理母線を $f_0$ とすると
$L\equiv s+4f_0$, $H\equiv s+6f_0$。
跡因子は $2L-K_T-2H\equiv2s+(2-2g_0)f_0$。有効性と $s^2=-4$ により
$s$ の係数はちょうど二である（一本引いても交点が負で、一般母線との交点は二）。

$s$ の入力次数二である。円錐曲線へ双有理ならその一般点で逆像は一つであり、
入力接方向と核の法方向が分離される。階数三で入力が同型になり跡の消滅に反する。
よって $s$ は直線 $\Lambda$ に二対一である。一般のその入力点では二枝とも滑らかである。
実際各枝で $s$ の方向は入力直線へ非分岐、核座標は一般法方向に非零で $s$ 上で定値、
かつ入力から正則に復元できるため入力の微分が単射になる。
他の正規化点はない。階数三では $M$ が可逆で評価も各枝でその引戻しを自明化するため、
跡合成の零因子は正規化の導手因子である。これはGorenstein超曲面への有限双対
$\mathcal Hom_D(f_*\mathcal O_T,\omega_D)$ を分数体内で書けば導手と $\omega_D$ の積になることによる。
一般横断平面での二滑らか枝の導手の各位数はその交点数である。
従って係数二から両枝の接平面は一致する。このことは直線の一般点の完備局所環を横断方向で
係数体上の平面曲線として見ても同じである。

共通接平面からの核写像の微分は $\Lambda$ の接方向を消し、両枝で同一の非零射影方向を与える。
これは頂点における錐母線の方向であり、$B_0$ の二点が同じ底パラメータ $t$ を持つことになる。
従って $s\to\Lambda$ は $B_0\to\mathbb P^1$ と射影同型の合成である。
$H$ と $g^*(E+3F)$ の数値差は底からの次数零の束であり、$s$ への制限が上の二重被覆の束として一致するので
$H=g^*(E+3F)$ が線形同値でも成立する。入力を跡分解すると
\[
          x=A_3(t)+z B_1(t)+ w C_{3-b}(t),\qquad w^2=P_{2b}(t)
\]
となる。$s$ 上の表示が底からの一次式 $B_1$（$\Lambda$ への同型）なので、$z$ の係数に反不変部分はない。
これは同じ射を与える二つの基点なしの座標列の比が $s$ 上の可逆な大域関数であることから従う。

各入力母線は $\Lambda$ に会う。$s$ の共役二点から出る二母線と $\Lambda$ が一平面に入る。
実際各線織枝の接平面は $\Lambda$ とその母線の張る平面であり、先に示した接平面の一致を使える。
従って $\Lambda$ のベクトル空間を法とする二列 $\bar A,\bar C$ は比例する。
$\bar C$ が非零なら、この共通方向は $\Lambda$ を通る平面束への写像を与え、その次数は $k\le3-b$。
定値なら入力曲面が一つの平面に入るので $k>0$ である。
一方その平面束の一般の一員は、八次曲面を $\Lambda$ と相異なる $2k$ 本の一般母線で切る。
$\Lambda$ の係数は二である。一般平面束の一員は $\Lambda$ の一般点で動く共通接平面と異なり、各滑らか枝を横断するからである。
残りは底の次数 $k$ の写像と二重被覆の一般非分岐繊維で、有限双有理入力射により各母線は一重である。
ゆえに $8=2+2k$、$k=3$ は上界に反する。
従って $C$ は $\Lambda$ 内にあり、非零である（零なら入力が二重被覆を経由する）。
最初の補題でこれも不可能である。
\end{proof}
\path{computations/verify_cofactor_cone_pencil.py} は平方根を変える局所計算と平面束の次数の矛盾を検算する。
